\begin{env}[5.10.1]
\label{IV.2.s5.10.1-fr}
Sean $X$ un preesquema localmente noetheriano y $\mathcal F$ un
$\mathcal O_X$-módulo coherente; para toda parte $T$ de $X$, pondremos
\begin{equation}
\operatorname{prof}_T(\mathcal F)
=\inf_{x\in T}\operatorname{prof}(\mathcal F_x)
\tag{5.10.1.1}\label{IV.2.s5.10.1.1-fr}
\end{equation}
\end{env}

\begin{proposition}[5.10.2]
\EGAReaderTarget{ega4.c04.prop.5-10-2}{IV.2.s5.10.2-fr}
Sean $X$ un preesquema localmente noetheriano, $Z$ una parte de $X$ estable
por especialización y $\mathcal F$ un $\mathcal O_X$-módulo cuasicoherente.
Las condiciones siguientes son equivalentes:
\begin{enumerate}
\item[\textup{a)}] $\mathcal F$ es $Z$-puro.
\item[\textup{b)}] $\operatorname{Ass}(\mathcal F)$ no corta $Z$.
\end{enumerate}
Si, además, $\mathcal F$ es coherente, estas condiciones son también
equivalentes a la siguiente:
\begin{enumerate}
\item[\textup{c)}] $\operatorname{prof}_Z(\mathcal F)\geq1$.
\end{enumerate}
\end{proposition}

Decir que $\mathcal F$ es $Z$-puro significa que, para todo abierto
noetheriano $V$ de $X$ y todo abierto $U\supset X\setmin Z$, el homomorfismo de
restricción
$\Gamma(V,\mathcal F)\longrightarrow\Gamma(V\cap U,\mathcal F)$
es inyectivo (5.9.8); pero, según (3.1.8), esto equivale a
$V\cap\operatorname{Ass}(\mathcal F)\subset U$, de donde se sigue la
equivalencia de \emph{a)} y \emph{b)}. Además, decir que
$x\in\operatorname{Ass}(\mathcal F)$ significa que ningún elemento de
$\mathfrak m_x$ es $\mathcal F_x$-regular (3.1.2); por tanto, cuando
$\mathcal F$ es coherente, esto también se escribe
$\operatorname{prof}(\mathcal F_x)=0$; en este caso se deduce inmediatamente
la equivalencia de \emph{b)} y \emph{c)}.

\begin{corollary}[5.10.3]
\label{IV.2.s5.10.3-fr}
Sea
$0\longrightarrow\mathcal F'\longrightarrow\mathcal F
\longrightarrow\mathcal F''\longrightarrow0$
una sucesión exacta de $\mathcal O_X$-módulos cuasicoherentes. Si
$\mathcal F$ es $Z$-puro, lo mismo ocurre con $\mathcal F'$; recíprocamente,
si $\mathcal F'$ y $\mathcal F''$ son $Z$-puros, lo mismo ocurre con
$\mathcal F$.
\end{corollary}

\oldpage[IV-2]{115}

Esto resulta de la forma (5.10.2, \emph{b)}) de la condición para que un
$\mathcal O_X$-módulo cuasicoherente sea $Z$-puro y de (3.1.7).

\begin{corollary}[5.10.4]
\label{IV.2.s5.10.4-fr}
Supongamos que $\mathcal F$ es coherente. Para que $\mathcal F$ sea $Z$-puro
en un punto $x\in X$, es necesario y suficiente que
$\operatorname{prof}_{Z_x}(\widetilde{\mathcal F}_x)\geq1$
\textup{(}con las notaciones de \textup{(5.9 6))}.
\end{corollary}

Esto resulta inmediatamente de (5.10.2) y (5.9.6).
